\section*{Capítulo \ref{ch:g11:prob} --- Probabilidade e variáveis aleatórias}

\begin{solution}{exo:g11:prob:1}
\omterm{def:g11:prob:model}{Modelo uniforme} sobre $52$ cartas. Copas:
$\frac{13}{52} = \frac14$. Reis: $\frac{4}{52} = \frac{1}{13}$. Copas ou
rei: $\frac{13 + 4 - 1}{52} = \frac{16}{52} = \frac{4}{13}$ (o rei de
copas é contado uma só vez). Nenhum dos dois:
$1 - \frac{4}{13} = \frac{9}{13}$.
\end{solution}

\begin{solution}{exo:g11:prob:2}
$S = 6$: os pares $(1,5), (2,4), (3,3), (4,2), (5,1)$, logo
$\P(S = 6) = \frac{5}{36}$.

$S \geq 10$: três pares para $10$, dois para $11$, um para $12$:
$\P(S \geq 10) = \frac{6}{36} = \frac16$.

$S$ par: as somas $2, 4, 6, 8, 10, 12$ têm
$1 + 3 + 5 + 5 + 3 + 1 = 18$ pares: \omterm{def:g10:proba:distribution}{probabilidade}
$\frac{18}{36} = \frac12$.
\end{solution}

\begin{solution}{exo:g11:prob:3}
As \omterm{def:g10:proba:distribution}{probabilidades} somam $1$: $p = 1 - 0.3 - 0.2 - 0.4 = 0.1$. Então
\[
\E(X) = 0.3 \times (-1) + 0.2 \times 0 + 0.4 \times 2 + 0.1 \times 5
= -0.3 + 0.8 + 0.5 = 1 .
\]
\end{solution}

\begin{solution}{exo:g11:prob:4}
$\E(Y) = 2 \times 3 - 5 = 1$; $\V(Y) = 2^2 \times 4 = 16$;
$\sigma(Y) = 4$.
\end{solution}

\begin{solution}{exo:g11:prob:5}
Os prêmios são $0$, $2$, $5$ euros com \omterm{def:g10:proba:distribution}{probabilidades} $0.5$, $0.3$,
$0.2$; subtraindo a aposta de $2$ euros, $G$ assume os valores
$-2, 0, 3$:
\begin{center}
\begin{tabular}{c|ccc}
$g$ & $-2$ & $0$ & $3$\\
\hline
$\P(G=g)$ & $0.5$ & $0.3$ & $0.2$
\end{tabular}
\end{center}
$\E(G) = -1 + 0 + 0.6 = -0.4$: o jogador perde $40$ centavos por partida
em \omterm{def:g11:stat:mean}{média} --- desfavorável.
\end{solution}

\begin{solution}{exo:g11:prob:6}
Retirando em ordem: $\P(X = 2) = \frac35 \times \frac24 = \frac{3}{10}$
(vermelha e depois vermelha);
$\P(X = 0) = \frac25 \times \frac14 = \frac{1}{10}$ (azul e depois
azul); logo $\P(X = 1) = 1 - \frac{3}{10} - \frac{1}{10} =
\frac{6}{10}$. Então
\[
\E(X) = 0 \times \frac1{10} + 1 \times \frac6{10} + 2 \times \frac3{10}
= \frac{12}{10} = 1.2 ,
\]
coerente com a intuição: em \omterm{def:g11:stat:mean}{média}, $\frac35$ das duas retiradas são
vermelhas.
\end{solution}

\begin{solution}{exo:g11:prob:7}
$\E(G^2) = 0.5 \times 4 + 0.3 \times 0 + 0.2 \times 9 = 3.8$, logo, pela
fórmula abreviada, $\V(G) = 3.8 - (-0.4)^2 = 3.8 - 0.16 = 3.64$ e
$\sigma(G) \approx 1.9$ euros.
\end{solution}

\begin{solution}{exo:g11:prob:8}
O resultado é $80$ menos o valor reclamado:
\begin{center}
\begin{tabular}{c|ccc}
$b$ & $80$ & $-920$ & $-4920$\\
\hline
$\P(B=b)$ & $0.94$ & $0.05$ & $0.01$
\end{tabular}
\end{center}
$\E(B) = 75.2 - 46 - 49.2 = -20$: a esse preço, a empresa \emph{perde}
$20$ euros por apólice em \omterm{def:g11:stat:mean}{média}. O prêmio é baixo demais para os riscos
cobertos (ele precisaria passar de $100$ euros para dar resultado
esperado positivo).
\end{solution}

\begin{solution}{exo:g11:prob:9}
A pontuação aleatória é $+1$ com \omterm{def:g10:proba:distribution}{probabilidade} $\frac14$ e $x$ com
\omterm{def:g10:proba:distribution}{probabilidade} $\frac34$:
\[
\E = \frac14 + \frac{3x}{4} = 0 \iff x = -\frac13 .
\]
Com a penalidade $-\frac13$, chutar ao acaso não traz nada em \omterm{def:g11:stat:mean}{média}.
\end{solution}

\begin{solution}{exo:g11:prob:10}
Duas moedas: elas coincidem (duas caras ou duas coroas) com
\omterm{def:g10:proba:distribution}{probabilidade} $\frac24 \times 2 = \frac12$. O ganho do jogador é $+1$ ou
$-1$ com \omterm{def:g10:proba:distribution}{probabilidades} iguais: \omterm{def:g11:prob:expectation}{esperança} $0$ para os dois jogadores ---
o jogo é justo.

Três moedas: as três coincidem (três caras ou três coroas) com
\omterm{def:g10:proba:distribution}{probabilidade} $\frac{2}{8} = \frac14$. O jogador recebe $2$ com
\omterm{def:g10:proba:distribution}{probabilidade} $\frac14$ e paga $1$ com \omterm{def:g10:proba:distribution}{probabilidade} $\frac34$:
$\E = \frac24 - \frac34 = -\frac14$. O jogo favorece Ana em $25$
centavos por rodada; seria justo se ela pagasse $3$ euros por uma trinca
coincidente.
\end{solution}

\begin{solution}{exo:g11:prob:11}
\emph{1.} O organizador recebe $1000 \times 5 = 5000$ euros e paga
$2000 + 5 \times 200 + 20 \times 25 = 3500$ euros ao todo: o lucro é a
\emph{constante} $1500$ (nenhuma aleatoriedade, uma vez vendidos todos os
bilhetes). O ganho $G$ de um único jogador (prêmio menos o bilhete de
$5$ euros):
\begin{center}
\begin{tabular}{c|cccc}
$g$ & $1995$ & $195$ & $20$ & $-5$\\
\hline\rule{0pt}{11pt}%
$\P(G=g)$ & $\dfrac{1}{1000}$ & $\dfrac{5}{1000}$ & $\dfrac{20}{1000}$
& $\dfrac{974}{1000}$
\end{tabular}
\end{center}

\emph{2.}
\[
\E(G) = \frac{1995 + 5 \times 195 + 20 \times 20 - 974 \times 5}{1000}
= \frac{1995 + 975 + 400 - 4870}{1000} = -1.5 .
\]
Cada jogador espera perder $1.50$ euro; em $1000$ jogadores isso dá
$1500$ euros --- exatamente o lucro do organizador. As contas fecham: o
que os jogadores perdem em \omterm{def:g11:stat:mean}{média} é precisamente o que o organizador
arrecada.
\end{solution}

\begin{solution}{pb:g11:prob:1}
\textbf{1.} $G$ assume $-2, -1, 0, 1, 2, 3$, cada um com
\omterm{def:g10:proba:distribution}{probabilidade} $\frac16$: $\E(G) = \frac{-2 - 1 + 0 + 1 + 2 +
3}{6} = \frac12$ euro.

\textbf{2.} $\E(G^2) = \frac{4 + 1 + 0 + 1 + 4 + 9}{6} =
\frac{19}{6}$; $V(G) = \frac{19}{6} - \frac14 = \frac{35}{12}
\approx 2.92$; $\sigma(G) \approx 1.71$.

\textbf{3.} $\E(2G - 1) = 2 \times \frac12 - 1 = 0$;
$V(2G - 1) = 4\,V(G) = \frac{35}{3}$.

\textbf{4.} $\E(\text{soma}) = 7$: preço justo $7$ euros. A $8$,
a perda esperada é de $1$ euro por partida.

\textbf{5.} Prêmio justo: $0.05 \times 800 = 40$ euros. Cobrando
$60$: lucro esperado de $20$ euros por apólice.

\textbf{6.} Dividir $5:3$ premia o \emph{passado} --- mas o bolo
pertence a quem ganhar o \emph{futuro}, e B ainda poderia
vencer. Dar tudo a A trata uma vitória provável como certa ---
a pequena chance de B vale alguma coisa.

\textbf{7.} Rodadas a jogar: $3$; resultados (vitória de A ou de
B por rodada): AAA, AAB, ABA, ABB, BAA, BAB, BBA, BBB. O
jogador B leva a série apenas ganhando as três: só BBB. A ganha
a série em $7$ dos $8$ casos: divisão justa de
$64 \times \frac78 = 56$ pistolas para A e $8$ para B.

\textbf{8.} Jogar as rodadas extras e inúteis não muda a vitória
de ninguém: uma vez que A tenha sua sexta vitória, as rodadas
seguintes são cerimônia. Mas fixar o horizonte em exatamente
$3$ rodadas torna os $8$ resultados igualmente prováveis ---
uma contagem honesta precisa de histórias do mesmo comprimento,
e preenchê-las não custa nada.

\textbf{9.} Em $5$--$5$: uma rodada decisiva, $32$ para cada um.
Em $5$--$4$: A recebe $\frac{64 + 32}{2} = 48$. Em $5$--$3$: a
rodada seguinte leva a $64$ (se A a vencer) ou à posição
$5$--$4$, que vale $48$: a parte de A é
$\frac{64 + 48}{2} = 56$ --- o $56 : 8$ de Fermat, rededuzido
caminhando de trás para a frente.

\textbf{10.} A precisa de $2$ e B de $3$: jogue $4$ rodadas
imaginárias ($16$ resultados). B leva a série com $3$ ou $4$
vitórias: $4 + 1 = 5$ resultados. Divisão: para B,
$64 \times \frac{5}{16} = 20$; para A, $44$.

\textbf{11.} O valor presente justo de um pagamento futuro
incerto é a sua \omterm{def:g11:prob:expectation}{esperança} --- a \omterm{def:g11:stat:mean}{média} ponderada pelas
\omterm{def:g10:proba:distribution}{probabilidades} do que pode acontecer. Os seguros (questão 5) e
as finanças (precificação de pagamentos incertos) são essa
frase transformada em indústria; os cassinos também, do outro
lado da mesa.

\textbf{12.} O bolo sempre vai parar em algum lugar: as duas
partes são variáveis aleatórias que somam a constante $64$,
logo, por linearidade, $\E_A + \E_B = 64$ --- as partes justas
esgotam o bolo, em qualquer placar.

\textbf{13.} Linearidade: $\E(X + Y) = 3.5 + 3.5 = 7$, sem
tabela --- verdadeiro até para dados colados um no outro.
\omterm{def:g11:prob:variance}{Variância} de um dado: $\E(X^2) - 3.5^2 = \frac{91}{6} -
\frac{49}{4} = \frac{35}{12}$; para dados independentes,
$V(X + Y) = \frac{35}{12} + \frac{35}{12} = \frac{35}{6}$.

\textbf{14.} Cada convidado fixado recupera o próprio chapéu com
\omterm{def:g10:proba:distribution}{probabilidade} $\frac1n$ (o chapéu dele é um entre $n$). Somando
sobre os $n$ convidados: $\E(X) = n \times \frac1n = 1$ --- um
convidado de sorte em \omterm{def:g11:stat:mean}{média}, tenha a festa dois chapéus ou dois
mil.

\textbf{15.} $n = 2$: as distribuições são ambos certos
($X = 2$) ou trocados ($X = 0$): $\E = 1$. $n = 3$: as seis
distribuições dão $X = 3, 1, 1, 1, 0, 0$:
$\E = \frac{6}{6} = 1$.

\textbf{16.} O cálculo somou as \omterm{def:g11:prob:expectation}{esperanças} dos $n$ indicadores
individuais (``o convidado $i$ pegou o próprio chapéu''), e a
linearidade soma \omterm{def:g11:prob:expectation}{esperanças} \emph{sem condições} --- a
dependência arruinaria um produto ou uma \omterm{def:g11:prob:variance}{variância}, mas nunca
uma soma de \omterm{def:g11:prob:expectation}{esperanças}. Esse é o superpoder da linearidade:
nenhuma pergunta sobre independência.

\textbf{17.} Os dois ganhos: $\E = 0.5 \times 10 - 5 = 0$ e
$0.005 \times 1000 - 5 = 0$. \omterm{def:g11:prob:variance}{Variâncias} (do ganho):
A: $\E(G^2) = 0.5 \times 25 + 0.5 \times 25 = 25$;
B: $0.005 \times 995^2 + 0.995 \times 25 \approx 4\,975$:
$\sigma_A = 5$ e $\sigma_B \approx 70.5$. As pessoas compram a
B: o que elas adquirem não é \omterm{def:g11:prob:expectation}{esperança} (nula nas duas), e sim a
\emph{dispersão} --- uma pequena chance de uma cauda direita
que muda a vida.

\textbf{18.} A primeira coroa no lançamento $n$ tem
\omterm{def:g10:proba:distribution}{probabilidade} $2^{-n}$ e paga $2^n$: cada $n$ contribui
$2^{-n} \times 2^n = 1$ para a \omterm{def:g11:prob:expectation}{esperança}, que é, portanto,
$1 + 1 + 1 + \dots$: infinita. E, no entanto, ninguém em sã
consciência paga $1\,000$ para jogar uma vez: com \omterm{def:g10:proba:distribution}{probabilidade}
$\frac{15}{16}$ o pagamento é no máximo $16$. A \omterm{def:g11:prob:expectation}{esperança}
resume o longo prazo de apostas \emph{repetíveis}; um jogo
único com uma cauda selvagem não se precifica pela \omterm{def:g11:stat:mean}{média}.

\textbf{19.} \omterm{def:g11:prob:expectation}{Esperanças}: com seguro, $-20$; sem seguro,
$-0.01 \times 1000 = -10$. A estratégia descoberta vence em
\omterm{def:g11:stat:mean}{média} --- para quem consegue absorver a perda e repetir a
aposta muitas vezes. Uma família que não sobrevive ao
$-1\,000$ não está no ramo das \omterm{def:g11:stat:mean}{médias}: ela paga $10$ a mais
para apagar a \omterm{def:g11:prob:variance}{variância}. As seguradoras tiram \omterm{def:g11:stat:mean}{médias}; as
famílias compram certeza.

\textbf{20.} Nascida da divisão de um bolo entre dois jogadores
impacientes; fortalecida pela linearidade, que soma \omterm{def:g11:prob:expectation}{esperanças}
sem pedir licença à dependência (os chapéus); cega ao risco,
que a \omterm{def:g11:prob:variance}{variância} --- as raspadinhas, o jogo de São Petersburgo, a
apólice da família --- tem de medir em seu lugar; e coroada no
ano que vem pela lei dos grandes números, que explica por que
as \omterm{def:g11:stat:mean}{médias} sempre acabam sendo cobradas.
\end{solution}
